%%=============================================================================
%% Inleiding
%%=============================================================================

\chapter{\IfLanguageName{dutch}{Inleiding}{Introduction}}%
\label{ch:inleiding}

\begin{comment}
    De inleiding moet de lezer net genoeg informatie verschaffen om het onderwerp te begrijpen en in te zien waarom de onderzoeksvraag de moeite waard is om te onderzoeken. In de inleiding ga je literatuurverwijzingen beperken, zodat de tekst vlot leesbaar blijft. Je kan de inleiding verder onderverdelen in secties als dit de tekst verduidelijkt. Zaken die aan bod kunnen komen in de inleiding
    
    \item context, achtergrond
    \item afbakenen van het onderwerp
    \item verantwoording van het onderwerp, methodologie
    \item probleemstelling
    \item onderzoeksdoelstelling
    \item onderzoeksvraag
    \item \ldots
\end{comment}

Men kan er niet aan onderuit, AI(Artificiële Intelligentie) is tegenwoordig overal. Het begint in dagdagelijks gebruikte software verwerkt te worden en gewild of ongewild komt men vroeg of laat in contact met AI in een of andere vorm. \\

Zelfs de meest sporadische gebruiker van digitale technologie zal nu en dan iets opzoeken met Google bijvoorbeeld en zullen dan merken dat het antwoord bovenaan de pagina een Google-Gemini gegenereerde response is, niet zomaar een zoekresultaat maar een output wat het resultaat is van de interne werking van een LLM(Large Language Model). \\

AI is overal, en nu is de vraag hoe dit kan ingezet worden op een veilige manier dat een positieve bijdrage levert in een Agile software development omgeving. Kan een AI agent helpen om bepaalde knelpunten in het development cyclus weg te werken en draagt het uiteindelijk bij tot een verhoogde efficiëntie?
% J j

\section{\IfLanguageName{dutch}{Probleemstelling}{Problem Statement}}%
\label{sec:probleemstelling}

\begin{comment}
    Uit je probleemstelling moet duidelijk zijn dat je onderzoek een meerwaarde heeft voor een concrete doelgroep. De doelgroep moet goed gedefinieerd en afgelijnd zijn. Doelgroepen als ``bedrijven,'' ``KMO's'', systeembeheerders, enz.~zijn nog te vaag. Als je een lijstje kan maken van de personen/organisaties die een meerwaarde zullen vinden in deze bachelorproef (dit is eigenlijk je steekproefkader), dan is dat een indicatie dat de doelgroep goed gedefinieerd is. Dit kan een enkel bedrijf zijn of zelfs één persoon (je co-promotor/opdrachtgever).
\end{comment}

In moderne softwareontwikkeling zijn APIs een cruciale schakel tussen systemen. Het ontwikkelen en testen van deze APIs vereist echter veel repetitieve en technische handelingen zoals het schrijven van testcases, het valideren van responses en het controleren van documentatie. In agile teams, waar snelheid en iteratie centraal staan, vormt dit een bottleneck. \\

Er is nood aan een intelligente assistent die deze taken kan ondersteunen of deels automatiseren. De vraag vanuit IT1(Webrand) rijst of een AI agent, gebaseerd op recente ontwikkelingen in natural language processing en machine learning, hierin een rol kan spelen.

%todo: mogelijks wat meer toevoegen of wat herschrijven / herschikken (nu is het copy-paste vanuit de inleiding van bp-voorstel)


\section{\IfLanguageName{dutch}{Onderzoeksvraag}{Research question}}%
\label{sec:onderzoeksvraag}

\begin{comment}
    Wees zo concreet mogelijk bij het formuleren van je onderzoeksvraag. Een onderzoeksvraag is trouwens iets waar nog niemand op dit moment een antwoord heeft (voor zover je kan nagaan). Het opzoeken van bestaande informatie (bv. ``welke tools bestaan er voor deze toepassing?'') is dus geen onderzoeksvraag. Je kan de onderzoeksvraag verder specifiëren in deelvragen. Bv.~als je onderzoek gaat over performantiemetingen, dan
\end{comment}

% J j
Dit onderzoek beschrijft het ondersteunen van software ontwikkeling met behulp van AI en de nieuwste ontwikkelingen binnen NLP(Natural Language Processing), de hoofdonderzoeksvraag luidt als volgt:

\begin{quote}
    \textit{Hoe kan een AI agent bijdragen aan het ondersteunen van het ontwikkel- en testproces van RESTful APIs binnen een agile softwareontwikkeltraject?}
\end{quote}

Hieronder is het probleem opgesplitst in een aantal deelvragen die zullen bijdragen tot een uitwerking van de hoofdonderzoeksvraag:

\begin{enumerate}
    \item Wat zijn de huidige uitdagingen bij het ontwikkelen en testen van RESTful APIs?
    \item Welke taken binnen dit proces zijn het meest repetitief of foutgevoelig?
    \item Welke bestaande tools en frameworks worden momenteel gebruikt voor API testing en documentatie?
    \item Welke AI-technieken (bv. LLMs, NLP, code generation) zijn geschikt om deze taken te ondersteunen?
    \item Hoe kan een AI agent geïntegreerd worden in een bestaande ontwikkelworkflow?
    \item Wat zijn de beperkingen en risico’s van het inzetten van een AI agent in dit domein?
\end{enumerate}


\section{\IfLanguageName{dutch}{Onderzoeksdoelstelling}{Research objective}}%
\label{sec:onderzoeksdoelstelling}
% J j
Het doel van dit onderzoek is om aan te tonen in hoeverre AI kan helpen bij het ontwikkelen en testen van APIs. Hiervoor zal een vergelijkende studie gedaan worden om te bepalen welke AI implementaties kunnen helpen om een proof of concept AI Agent te ontwikkelen. Er zal onderzocht worden hoe deze AI agent kan ingezet worden voor het ondersteunen van de API ontwikkelingsprocessen. \\

Kan de Agent ingezet worden voor belangrijke taken zoals het automatisch genereren van testcases, analyse van foutmeldingen en het aanvullen of genereren van documentatie? \\

Het onderzoek zal aanwijzen of het inzetten van AI agents de efficiëntie van het API development process al dan niet kan optimaliseren.

%todo: consider "...hoe deze AI agent kan ingezet worden..." -> hoe deze AI agent(s) kunnen ingezet worden, multiple agents?

\begin{comment}
    Wat is het beoogde resultaat van je bachelorproef? Wat zijn de criteria voor succes? Beschrijf die zo concreet mogelijk. Gaat het bv.\ om een proof-of-concept, een prototype, een verslag met aanbevelingen, een vergelijkende studie, enz.
\end{comment}

\section{\IfLanguageName{dutch}{Opzet van deze bachelorproef}{Structure of this bachelor thesis}}%
\label{sec:opzet-bachelorproef}

% Het is gebruikelijk aan het einde van de inleiding een overzicht te
% geven van de opbouw van de rest van de tekst. Deze sectie bevat al een aanzet
% die je kan aanvullen/aanpassen in functie van je eigen tekst.

De rest van deze bachelorproef is als volgt opgebouwd:

In Hoofdstuk~\ref{ch:stand-van-zaken} wordt een overzicht gegeven van de stand van zaken binnen het onderzoeksdomein, op basis van een literatuurstudie.

In Hoofdstuk~\ref{ch:methodologie} wordt de methodologie toegelicht en worden de gebruikte onderzoekstechnieken besproken om een antwoord te kunnen formuleren op de onderzoeksvragen.

% TODO: Vul hier aan voor je eigen hoofstukken, één of twee zinnen per hoofdstuk

In Hoofdstuk~\ref{ch:conclusie}, tenslotte, wordt de conclusie gegeven en een antwoord geformuleerd op de onderzoeksvragen. Daarbij wordt ook een aanzet gegeven voor toekomstig onderzoek binnen dit domein.